\documentclass[11pt,a4paper]{article}
\usepackage[utf8]{inputenc}
\usepackage[T1]{fontenc}
\usepackage{lmodern}
\usepackage[brazil]{babel}
\usepackage{amsmath,amssymb}
\usepackage{booktabs,tabularx,array,xltabular}
\usepackage{enumitem}
\usepackage{tikz}
\usetikzlibrary{arrows.meta,positioning}
\usepackage[a4paper,top=2.1cm,bottom=2.1cm,left=2.3cm,right=2.3cm]{geometry}
\usepackage{microtype}
\usepackage[hidelinks]{hyperref}

\setlist{itemsep=1pt,topsep=3pt,parsep=0pt}
\setlength{\parskip}{3pt}
\setlength{\parindent}{0pt}
\renewcommand{\arraystretch}{1.15}
\newcommand{\mods}[1]{\ \mathrm{mod}\ #1}
\newcolumntype{Y}{>{\raggedright\arraybackslash}X}

\begin{document}

% ---------- Cabeçalho ----------
\begin{center}
{\fontfamily{ptm}\selectfont
{\Large\bfseries Projeto Final\\[2pt] Sistema de Votação Eletrônica Seguro\par}
\vspace{14pt}
{\large\bfseries Alan dos Santos Dias, 232007830\\
Enzo Cardoso Martins, 242010033\\
Felipe Lauterjung Caselli, 241032401\\
Leandro Souza da Silva, 211031109\\
Rafael Medeiros de Alencar, 241020760\\
\par}
\vspace{12pt}
{\large $^{1}$Departamento de Ciência da Computação -- Universidade de Brasília (UnB)\\
Segurança Computacional\par}
}
\end{center}
\vspace{4pt}

%==============================================================
\section{Introdução e Visão Geral}
Este trabalho projeta e implementa, em Python, um sistema de votação eletrônica que conduz uma eleição completa (configuração, cadastro, votação, encerramento, apuração e auditoria) sem exigir confiança em nenhuma entidade isolada, nem no administrador. A tese central é que a segurança vem da matemática e não de uma autoridade central: o voto é cifrado com \emph{ElGamal exponencial}, a chave de decifração é dividida entre $N\ge 3$ autoridades e nunca existe inteira em lugar algum, as cédulas ficam num \emph{quadro público} encadeado por SHA-256 e assinadas com RSA, e um auditor independente verifica tudo a partir de dados públicos.

A eleição simulada tem 3 candidatos (Alice, Bob, Carol), 3 autoridades e 5 eleitores cadastrados. Só usamos primitivas básicas: aritmética modular de Python (\texttt{pow}), \texttt{hashlib}, \texttt{secrets} e a biblioteca \texttt{cryptography} apenas para RSA. O ElGamal, a chave conjunta, a apuração homomórfica e o quadro público foram implementados por nós, como exige o enunciado.

%==============================================================
\section{Arquitetura e Responsabilidades}
\subsection{Visão geral}
O sistema tem cinco entidades (Figura~\ref{fig:arq}). Nenhuma delas, sozinha, consegue ler um voto individual nem alterar o resultado sem ser detectada.

\begin{figure}[h]
\centering
\begin{tikzpicture}[node distance=2.4cm and 3.1cm,
  box/.style={draw,rounded corners,align=center,minimum width=2.5cm,minimum height=0.9cm,font=\small},
  arr/.style={-{Stealth[length=2mm]},thick},
  lab/.style={font=\scriptsize,align=center}]
 \node[box] (el) {Eleitor\\(cliente)};
 \node[box,right=of el] (sv) {Servidor\\de votação};
 \node[box,right=of sv] (au) {Autoridades\\($N\ge 3$)};
 \node[box,below=of sv] (qp) {Quadro\\Público};
 \node[box,right=of qp] (ad) {Auditor\\independente};
 \draw[arr] (el) -- node[lab,above]{cédula cifrada\\+ assinatura RSA} (sv);
 \draw[arr,dashed] (sv.south west) -- node[lab,below left]{recibo\\assinado} (el.south east);
 \draw[arr] (sv) -- node[lab,above]{$C_1$ \,/\, $h_i$} (au);
 \draw[arr,dashed] (au.south west) -- node[lab,below right=-1pt]{$d_i = C_1^{x_i}$} (sv.south east);
 \draw[arr] (sv) -- node[lab,right]{append} (qp);
 \draw[arr] (qp) -- node[lab,above]{lê (somente\\dados públicos)} (ad);
\end{tikzpicture}
\caption{Entidades e fluxo de mensagens.}\label{fig:arq}
\end{figure}

\subsection{Eleitores (cliente) --- \texttt{votacao.py}}
Cada eleitor possui um par RSA cadastrado antes da votação. O cliente monta o \emph{vetor de votos} (1 na posição do candidato escolhido, 0 nas demais), cifra cada posição com um $r$ novo, serializa $\{\text{id},\text{cédula}\}$ em JSON canônico, assina com RSA-PSS e envia o pacote. Ao receber o recibo, confere a assinatura do servidor.

\subsection{Servidor de votação (administrador) --- \texttt{setup\_eleicao.py}, \texttt{votacao.py}}
No setup, recebe as chaves parciais $h_i$ e calcula $h=\prod h_i \bmod p$. Guarda apenas as chaves \emph{públicas} RSA dos eleitores. Na votação, valida cada pacote (formato, elegibilidade, unicidade, assinatura), anexa o bloco ao quadro e emite o recibo assinado. Na apuração, multiplica as cifras e combina as decifrações parciais. Ele \emph{não} possui nenhuma parcela de $x$ e, portanto, não decifra nada.

\subsection{Autoridades eleitorais --- \texttt{setup\_eleicao.py}, \texttt{apuracao.py}}
São $N\ge 3$ entidades independentes. Cada uma gera $x_i$ (que nunca sai de sua posse), publica $h_i=g^{x_i}$ e, na apuração, devolve $d_i=C_1^{x_i}$ para cada candidato.

\subsection{Quadro público --- \texttt{votacao.py}}
Arquivo JSON somente de acréscimo. Contém a configuração pública da eleição ($p$, $g$, $h$, os $h_i$, candidatos, chave pública do servidor, lista de eleitores com suas chaves públicas), o \texttt{hash\_genesis} e a lista de blocos (id do eleitor, cédula cifrada, assinatura, instante, hash anterior e hash atual).

\subsection{Auditor independente --- \texttt{auditor.py}}
Programa separado do servidor, que lê \emph{somente} \texttt{quadro\_publico.json}. Não acessa segredos nem o estado interno do servidor (detalhes na Seção~\ref{sec:aud}).

%==============================================================
\section{Modelo de Ameaças}
\begin{xltabular}{\textwidth}{@{}>{\bfseries\raggedright\arraybackslash}p{2.9cm}YY@{}}
\toprule
Adversário & O que consegue tentar & Contramedida e limite\\
\midrule
Administrador / servidor malicioso &
Forjar cédulas; alterar, remover ou reordenar cédulas; votar por eleitores que não votaram; decifrar votos; recusar cédulas. &
Forjar exige a chave RSA do eleitor. Alterar/reordenar quebra o encadeamento e a assinatura. Não tem $x_i$, logo não decifra. \emph{Limite:} pode censurar uma cédula antes de aceitá-la (o eleitor percebe por não receber recibo) e pode cortar as \emph{últimas} cédulas sem quebrar a cadeia (só quem tem recibo percebe, pois seu \texttt{hash\_bloco} sumiria do quadro).\\
Autoridade desonesta &
Decifrar sozinha; recusar-se a participar; enviar $d_i$ falso. &
Um único $x_i$ não decifra nada. Recusa ou $d_i$ falso fazem a apuração falhar (testado na demo), sem produzir resultado errado. \emph{Limite:} a falha não aponta o culpado (não há prova de decifração correta) e o esquema é $N$-de-$N$ (Seção~\ref{sec:iii}).\\
Eleitor desonesto &
Votar duas vezes; votar por outro; enviar cédula malformada ou com valor $\ne 0,1$; provar o voto a terceiros. &
Unicidade e assinatura RSA barram os dois primeiros; checagem de formato barra a malformada. \emph{Limite:} sem prova de conhecimento zero, valores $\ne0,1$ não são detectados (Seção~\ref{sec:i}); recibo (Seção~\ref{sec:iv}).\\
Observador da rede &
Ler pacotes; modificar ou reenviar cédulas. &
Só vê cifras probabilísticas (nada sobre o voto). Modificar invalida a assinatura; reenviar esbarra na unicidade. \emph{Limite:} a simulação roda em um processo, sem TLS; quem votou (o id) é público, apenas o conteúdo é secreto.\\
\bottomrule
\end{xltabular}

%==============================================================
\section{Cifragem dos Votos e Apuração Homomórfica}
\subsection{ElGamal exponencial vetorial}
Sejam $p$ primo, $g$ gerador e $h$ a chave pública conjunta. Para cada candidato $k$, o eleitor escolhe $m_k\in\{0,1\}$ (1 só para o escolhido), sorteia $r_k$ com \texttt{secrets} e produz
\[
(c_{1,k},c_{2,k})=\bigl(g^{r_k}\bmod p,\; g^{m_k}\,h^{r_k}\bmod p\bigr).
\]
Como $r_k$ é novo a cada cifra, dois eleitores com o mesmo voto geram cifras completamente distintas, e as posições 0 e 1 da mesma cédula são indistinguíveis para quem não conhece $x$ (segurança semântica sob DDH). O vetor binário, em vez do id do candidato, é o que permite somar por candidato.

\subsection{Apuração sem decifrar votos individuais}
Para cada candidato, o servidor multiplica as cifras de todas as $n$ cédulas:
\[
C_1=\prod_{j=1}^{n}c_{1,j}=g^{R},\qquad
C_2=\prod_{j=1}^{n}c_{2,j}=g^{M}\,h^{R},\qquad R=\sum_j r_j,\;\; M=\sum_j m_j .
\]
$M$ é o total de votos do candidato, e $h^R$ é a máscara que impede ler $g^M$. Nenhuma cédula individual é decifrada: só o agregado.

%==============================================================
\section{Chave Conjunta e Decifração Distribuída}
\subsection{Geração cooperativa da chave}
Cada autoridade $i$ sorteia $x_i$ e publica $h_i=g^{x_i}\bmod p$. O servidor calcula
\[
h=\prod_{i=1}^{N}h_i=g^{\sum_i x_i}\bmod p ,
\]
ou seja, a chave privada conjunta é $x=\sum_i x_i$, mas \emph{ninguém a conhece} (cada $x_i$ fica só com sua autoridade). Não há chave mestra com o administrador.

\subsection{Decifração parcial}
O servidor envia o valor público $C_1$ de cada candidato. A autoridade $i$ devolve $d_i=C_1^{x_i}\bmod p$. Obter $x_i$ a partir de $(C_1,d_i)$ equivale a calcular um logaritmo discreto em um grupo de 2048 bits, inviável.

\subsection{Recombinação e extração}
O servidor calcula
\[
S=\prod_{i=1}^{N}d_i=C_1^{\sum x_i}=g^{R x}=h^{R},\qquad
C_2\cdot S^{-1}\bmod p=g^{M}.
\]
Para obter $M$ testamos $g^0,g^1,\dots,g^n$ até coincidir (\texttt{decifrar\_logaritmo\_discreto}). A busca exaustiva é viável porque $M\le n$ (número de eleitores), e inofensiva porque só se aplica ao total, nunca a uma cédula.

%==============================================================
\section{Integridade do Quadro e Autenticação dos Eleitores}
\subsection{Estrutura \emph{append-only} e hash encadeado}
O quadro nasce com $H_0=\text{SHA-256}(\text{config. canônica})$, que o amarra aos parâmetros da eleição. Cada bloco $k$ (dados $D_k$: índice, id, chave pública, cédula, assinatura, instante) recebe
\[
H_k=\text{SHA-256}\bigl(H_{k-1}\,\Vert\,\text{canônico}(D_k)\bigr),
\]
e guarda $H_{k-1}$ e $H_k$. Alterar, remover ou reordenar qualquer bloco muda seu hash e invalida todos os seguintes. A única operação de escrita do quadro é \texttt{adicionar\_bloco}, que apenas acrescenta no fim.

\subsection{Serialização canônica}
Hash e assinatura dependem de bytes idênticos em quem assina e em quem verifica. Usamos \texttt{json.dumps(sort\_keys=True, separators=(',',':'))}: sem isso, a ordem das chaves do dicionário ou espaços diferentes quebrariam hashes e assinaturas. Os números grandes são mantidos como inteiros, para que o mesmo valor gere sempre o mesmo texto.

\subsection{Autenticação e não-repúdio}
O eleitor assina $\{\text{id},\text{cédula}\}$ com RSA-2048, padding PSS e SHA-256. A assinatura vai no quadro, portanto qualquer pessoa verifica, com a chave pública registrada na configuração, que aquela cédula foi emitida por aquele eleitor e não foi alterada; o eleitor não pode negar tê-la enviado. O servidor assina, com sua própria chave RSA, o recibo de cada voto.

%==============================================================
\section{Auditoria e Verificação Individual}\label{sec:aud}
\subsection{Verificação individual (recibo)}
O recibo contém $\{\text{id},\text{índice},\text{hash\_bloco},\text{instante}\}$ assinado pelo servidor. O eleitor confere (comando \texttt{verificar-recibo}) que (i) a assinatura do servidor é válida e (ii) o bloco de mesmo índice no quadro público tem o mesmo \texttt{hash\_bloco} e o seu id. Se o servidor tivesse descartado ou trocado a cédula, o recibo não bateria com o quadro.

\subsection{Auditor independente universal}
O auditor lê só \texttt{quadro\_publico.json} e executa: (1) recálculo do \texttt{hash\_genesis}; (2) verificação da cadeia de hashes do primeiro ao último bloco; (3) verificação da assinatura RSA de cada eleitor sobre sua cédula; (4) elegibilidade (todo id consta na lista cadastrada); (5) unicidade (nenhum id repetido); (6) formato das cédulas (uma cifra por candidato, valores em $[1,p-1]$); e (7) recálculo independente do produto homomórfico $(C_1,C_2)$ de cada candidato. A \emph{independência} é de dados e de privilégios: não recebe segredos, não consulta o estado do servidor e não decifra nada, embora reutilize funções utilitárias do repositório (hash, verificação RSA, serialização). Na demonstração, alterar uma única cifra de um bloco faz falharem a cadeia de hashes e a assinatura do eleitor, e o auditor acusa o bloco exato.

%==============================================================
\section{Parâmetros e Algoritmos Criptográficos}
\begin{tabularx}{\textwidth}{@{}>{\bfseries\raggedright\arraybackslash}p{2.7cm}p{4.2cm}Y@{}}
\toprule
Componente & Escolha & Justificativa\\
\midrule
Grupo ElGamal & RFC 3526 grupo 14: $p$ primo seguro de 2048 bits, $g=2$ &
$p=2q+1$ com $q$ primo: não há subgrupos pequenos (evita Pohlig--Hellman). Como $p\equiv7\pmod 8$, $2$ é resíduo quadrático e gera o subgrupo de ordem $q$. 2048 bits dão cerca de 112 bits de segurança (NIST SP 800-57). Grupo padronizado é permitido pelo enunciado e evita erros de geração própria.\\
Assinaturas & RSA-2048, $e=65537$, PSS (MGF1/SHA-256, sal máximo) &
PSS é probabilístico e tem prova de segurança no modelo do oráculo aleatório, ao contrário do PKCS\#1 v1.5. Usado por eleitores e servidor.\\
Hash & SHA-256 &
Resistente a colisão e pré-imagem; padrão para encadeamento e para a assinatura PSS.\\
Aleatoriedade & \texttt{secrets} (entropia do SO) &
$x_i$ e $r$ precisam ser imprevisíveis; \texttt{random} (Mersenne Twister) permitiria reconstruir o estado e prever $r$.\\
Aritmética & \texttt{pow(b,e,m)}, \texttt{pow(a,-1,m)} &
Exponenciação modular por quadrados e multiplicações, sem estourar memória; o inverso usa Euclides estendido.\\
\bottomrule
\end{tabularx}

%==============================================================
\section{Análise de Segurança}
\subsection{Requisitos mínimos}
\begin{tabularx}{\textwidth}{@{}>{\bfseries\raggedright\arraybackslash}p{3.1cm}Y@{}}
\toprule
Requisito & Como é atendido (e limite)\\
\midrule
Sigilo do voto & Cifra ElGamal probabilística; sem $x$ ninguém decifra; ordem e instante não revelam nada pois as cifras são independentes do voto. \emph{Limite:} conluio de todas as autoridades (Seção~\ref{sec:ii}).\\
Integridade & Hash encadeado e assinaturas RSA. \emph{Limite:} corte das últimas cédulas só é notado por quem tem recibo.\\
Elegibilidade & Chaves públicas homologadas na configuração; servidor e auditor conferem.\\
Unicidade & Servidor recusa id que já votou; auditor confere ids repetidos.\\
Justiça & Resultado parcial exige $d_i$ de \emph{todas} as autoridades, e a CLI só apura após \texttt{encerrar}. \emph{Limite:} depende de as autoridades não fornecerem $d_i$ antes do fim.\\
Verificabilidade individual & Recibo assinado, conferível contra o quadro.\\
Auditabilidade & Auditor independente, só com dados públicos (Seção~\ref{sec:aud}).\\
\bottomrule
\end{tabularx}

\subsection{(i) Eleitor cifra um valor diferente de 0 ou 1}\label{sec:i}
O sistema \emph{não} impede isso: o servidor só vê $(c_1,c_2)$ e não consegue saber o $m$ dentro da cifra, e não há prova de conhecimento zero (ZKP) de que $m\in\{0,1\}$. Se um eleitor cifra $m=10$ (ou $m=-1$, ou vários 1 na mesma cédula), o total do candidato fica corrompido. Se a soma passa do número de cédulas, a busca exaustiva falha e a apuração aborta (um único eleitor desonesto consegue travar a contagem); se a soma ainda cabe na faixa, a fraude \emph{passa despercebida}. É a principal limitação da solução. A correção seria exigir, em cada posição, uma prova disjuntiva de Chaum--Pedersen de que $m\in\{0,1\}$, e uma prova de que a soma da cédula é 1; não a implementamos.

\subsection{(ii) Todas as autoridades conspiram}\label{sec:ii}
O sigilo cai. Elas somam os segredos, $x=\sum_i x_i$, e decifram \emph{cada cédula} do quadro, que é público: $g^{m}=c_2\cdot(c_1^{x})^{-1}$, com $m\in\{0,1\}$ trivial de ler. Isso mostra o preço do modelo: o sigilo exige que \emph{pelo menos uma} autoridade seja honesta. Um subconjunto estrito não basta, pois falta $c_1^{x_k}$ da autoridade excluída.

\subsection{(iii) Uma autoridade não participa da apuração}\label{sec:iii}
Sem $d_k$, o servidor obtém $S'=\prod_{i\ne k}d_i=h^R/C_1^{x_k}$ e, portanto,
$C_2\cdot S'^{-1}=g^{M}\,C_1^{x_k}=g^{M+Rx_k}$, que parece um elemento aleatório do grupo. A chance de coincidir com algum $g^0,\dots,g^n$ é desprezível, então a busca falha e \emph{nenhum resultado é divulgado} (demonstrado na CLI com \texttt{apurar --sem-autoridade 3}). O custo é de disponibilidade: o esquema é $N$-de-$N$, e uma só autoridade ausente bloqueia a eleição. Um esquema limiar (compartilhamento de Shamir) resolveria, ao custo de maior complexidade.

\subsection{(iv) O recibo permite provar a terceiros em quem se votou?}\label{sec:iv}
Não por si só. O recibo contém apenas $\{\text{id},\text{índice},\text{hash\_bloco},\text{instante}\}$, assinado pelo servidor; não traz o voto, a cédula nem o $r$. Ele prova \emph{que} a cédula foi registrada, não \emph{em quem} se votou, o que dificulta compra de votos e coerção. Há uma ressalva honesta: o cliente gera $r$ e o descarta, mas um cliente modificado poderia guardá-lo. Com $(m,r)$ qualquer um verifica $c_1=g^r$ e $c_2\,g^{-m}=h^r$, e o voto é provado. Logo o sistema não é \emph{à prova de coerção}; isso exigiria, por exemplo, recifragem ou voto revogável.

\subsection{Outras limitações}
(a) Não há prova de conhecimento de $x_i$: a última autoridade a publicar poderia escolher $h_N=Y\cdot(\prod_{i<N}h_i)^{-1}$ e passar a ser a única a conhecer a chave conjunta; mitigaríamos com provas de Schnorr ou \emph{commit--reveal} no setup. (b) Não há prova de que $d_i$ foi calculado corretamente (Chaum--Pedersen), então uma autoridade que sabota só é detectada pela falha, sem ser identificada. (c) O cadastro gera as chaves dos eleitores (simulação); na prática cada eleitor geraria a sua. (d) O resultado final não é publicado com provas, portanto o auditor valida $C_1,C_2$, mas não confirma $M$ sozinho.

%==============================================================
\section{Dificuldades e Aprendizados}
\textbf{Aritmética com inteiros de 2048 bits.} Entender por que \texttt{pow(b,e,m)} resolve o que \texttt{(b**e)\%m} não resolve (memória) e usar o inverso modular em vez de ``dividir''. Aprendizado: a divisão no grupo é multiplicar por $a^{-1}$.

\textbf{Consistência da chave conjunta.} Garantir que $h=g^{\sum x_i}$ de fato bata com os segredos distintos de cada autoridade, e que nenhuma delas reaparecesse duas vezes. Aprendizado: a descentralização da confiança é só aritmética de expoentes.

\textbf{Serialização canônica.} Hash e assinatura quebravam quando a ordem das chaves ou o tipo dos números (inteiro $\times$ texto) mudava entre quem assina e quem verifica. Aprendizado: segurança criptográfica depende de bytes idênticos.

\textbf{Reconstrução da máscara e depuração do logaritmo discreto.} Quando $S$ está errado, o resultado não é ``um pouco errado'', é aleatório e a busca simplesmente não acha. Aprendizado: foi esse comportamento que nos deu a resposta da pergunta (iii).

\textbf{Auditor desacoplado.} Fazer o auditor funcionar só com o arquivo público, sem consultar o servidor. Aprendizado: o que se quer verificável precisa estar no quadro (por isso a configuração com as chaves dos eleitores está no \texttt{hash\_genesis}).

%==============================================================
\section{Instruções de Compilação e Execução}
Requer Python 3.9+ (Linux/Ubuntu). A partir da raiz do projeto:
\begin{verbatim}
python3 -m venv venv && source venv/bin/activate
pip install -r requirements.txt          # cryptography
python3 src/cli.py demo                  # roteiro completo da apresentação
\end{verbatim}
O \texttt{demo} executa: eleição com 3 autoridades e 5 eleitores; voto duplo rejeitado; adulteração de uma cédula detectada pelo auditor; apuração sem uma autoridade (falha); apuração final. Comandos individuais (\texttt{setup}, \texttt{votar}, \texttt{verificar-recibo}, \texttt{encerrar}, \texttt{adulterar}, \texttt{auditar}, \texttt{apurar [--sem-autoridade N]}, \texttt{status}) estão no \texttt{README.md}. O auditor também roda sozinho: \texttt{python3 src/auditor.py eleicao/publico/quadro\_publico.json}.

\end{document}
